\section*{Hoofdstuk \ref{ch:b1:carbohydrates} --- Koolhydraten}

\begin{solution}{exo:b1:carbohydrates:1}
Glucose: hexose, \omterm{def:b1:carbohydrates:monosaccharide}{aldose}. Fructose: hexose, \omterm{def:b1:carbohydrates:monosaccharide}{ketose}. Ribose: pentose, \omterm{def:b1:carbohydrates:monosaccharide}{aldose}.
Glyceraldehyde: triose, \omterm{def:b1:carbohydrates:monosaccharide}{aldose}.
\end{solution}

\begin{solution}{exo:b1:carbohydrates:2}
De carbonylkoolstof na sluiting van de ring, die de ringzuurstof en een
hydroxylgroep draagt en die weer tot het aldehyde of keton kan opengaan. In
sucrose zitten beide anomere koolstoffen vast in de \omterm{def:b1:carbohydrates:glycosidic}{glycosidebinding}; geen van
beide kan opengaan, zodat er geen aldehyde ontstaat en het koper niet wordt
gereduceerd.
\end{solution}

\begin{solution}{exo:b1:carbohydrates:3}
\omterm{def:b1:carbohydrates:polysaccharide}{Zetmeel}: glucose, $\alpha(1{\to}4)$ met $\alpha(1{\to}6)$-vertakkingen
(amylopectine). \omterm{def:b1:carbohydrates:polysaccharide}{Glycogeen}: hetzelfde, sterker vertakt. \omterm{def:b1:carbohydrates:polysaccharide}{Cellulose}: glucose,
$\beta(1{\to}4)$. \omterm{def:b1:carbohydrates:polysaccharide}{Chitine}: N-acetylglucosamine, $\beta(1{\to}4)$.
\end{solution}

\begin{solution}{exo:b1:carbohydrates:4}
C1 omlaag, C2 omlaag, C3 omhoog, C4 omlaag. $\alpha$-D-galactose is hetzelfde met
de C4-hydroxylgroep omhoog.
\end{solution}

\begin{solution}{exo:b1:carbohydrates:5}
Bij de $\alpha$-binding wijst de anomere hydroxylgroep omlaag, en kan elk residu
onder dezelfde hoek als het vorige worden aangehecht: de keten draait gestaag door
en sluit zich tot een helix. Bij de $\beta$-binding wijst zij omhoog, en kan een
residu alleen omgeklapt worden aangehecht: de keten loopt recht. Jodiummoleculen
passen in de amylosehelix en hun elektronentoestand verandert (blauw); een vlak
lint biedt geen holte.
\end{solution}

\begin{solution}{exo:b1:carbohydrates:6}
Jodium: alleen \omterm{def:b1:carbohydrates:polysaccharide}{zetmeel} wordt blauwzwart. Benedict op de andere twee: alleen
glucose geeft een rood neerslag; sucrose blijft blauw. Bevestig sucrose door haar
met verdund \omterm{def:b1:water-small-molecules:ph}{zuur} te verwarmen, te neutraliseren en de proef van Benedict te
herhalen: nu positief.
\end{solution}

\begin{solution}{exo:b1:carbohydrates:7}
Ketens van 12 residuen: ongeveer \num{4000} ketens, waarvan de helft buitenste:
zo'n \num{2000} uiteinden. Amylose: één uiteinde. Om \num{5000} glucosen vrij te
geven bij één per seconde per uiteinde: \omterm{def:b1:carbohydrates:polysaccharide}{glycogeen} \qty{2.5}{s}; amylose
\qty{5000}{s}, bijna anderhalf uur.
\end{solution}

\begin{solution}{exo:b1:carbohydrates:8}
$\sigma = 0.8\times 25/(2\times 0.5) = \qty{20}{MPa}$: een vijftigste van de
sterkte van de fibrillen; de wand heeft een ruime veiligheidsmarge, en die is
nodig omdat de fibrillen maar een kwart van de wand uitmaken.
\end{solution}

\begin{solution}{exo:b1:carbohydrates:9}
Lactase verbreekt de $\beta(1{\to}4)$-binding tussen galactose en glucose; zonder
dat enzym wordt lactose niet opgenomen (alleen \omterm{def:b1:carbohydrates:monosaccharide}{monosachariden} passeren het
\omterm{def:b1:body-plans-tissues:epithelium}{epitheel}). In de dikke darm vergisten bacteriën haar tot zuren en gas (krampen,
opgeblazen gevoel), en de onverteerde suiker en de zuren trekken osmotisch water
het lumen in: diarree.
\end{solution}

\begin{solution}{exo:b1:carbohydrates:10}
Bacteriën en protisten in de pens scheiden cellulasen af en vergisten de glucose
tot korte \omterm{def:b1:lipids:fattyacid}{vetzuren}, die de koe opneemt en oxideert; de microben zelf worden later
verteerd. De vergisting moet aan de dunne darm voorafgaan, zodat haar producten
het opnemende oppervlak bereiken en het microbiële eiwit in de lebmaag en de darm
terechtkomt. \omterm{def:b1:carbohydrates:polysaccharide}{Cellulose}: \qty{4}{kg}; vergist \qty{2.4}{kg}; \qty{29}{MJ} per dag,
het grootste deel van de energie van de koe.
\end{solution}

\begin{solution}{exo:b1:carbohydrates:11}
\omterm{def:b1:cell-unit-of-life:cell}{O-cellen} dragen geen van beide antigenen en lokken bij niemand een antistof uit;
\omterm{def:b1:mammal-organization:compartments}{O-plasma} bevat echter antistoffen tegen zowel A als B, zodat iemand met O elke
ontvangen A- of \omterm{def:b1:cell-unit-of-life:cell}{B-cel} aanvalt. Het afweersysteem van iemand met A heeft B nooit
gezien en maakt anti-B-antistoffen (tegen darmbacteriën met vergelijkbare
suikers), maar verdraagt A: \omterm{def:b1:mammal-organization:compartments}{A-plasma} bevat alleen anti-B.
\end{solution}

\begin{solution}{exo:b1:carbohydrates:12}
Goed gezien: glucose is de brandstof van vrijwel elke \omterm{def:b1:cell-unit-of-life:cell}{cel} en \omterm{def:b1:carbohydrates:polysaccharide}{cellulose} het meest
voorkomende organische molecuul op aarde, en beide zijn polymeren van, of
hetzelfde als, D-glucose. Eén binding verandert alles: $\alpha$ geeft een spiraal
die enzymen openen en \omterm{def:b1:cell-unit-of-life:cell}{cellen} verbranden; $\beta$ geeft een vezel die enzymen,
water en trek weerstaat en die de meeste dieren helemaal niet kunnen gebruiken.
Biologie is niet de lijst van moleculen maar de meetkunde van hun verbinding ---
dezelfde chemie kan voedsel of geraamte zijn naar gelang één stereochemische
keuze, en \omterm{def:b1:organism-environment:organism}{organismen} worden gescheiden door de vraag welke enzymen zij bezitten
om haar ongedaan te maken.
\end{solution}

\begin{solution}{pb:b1:carbohydrates:1}
\textbf{1.} $2^{t-1}$: 1, 2, 4, en $2^{11} = 2048$.
\textbf{2.} $2^{12} - 1 = 4095$.
\textbf{3.} $4095\times 13 = \num{53235}$, ongeveer \num{53000}.
\textbf{4.} $\num{53235}\times 162 = \qty{8.6e6}{Da}$;
$\qty{1.43e-17}{g}$.
\textbf{5.} $2048/4095$: de helft. \num{2048} niet-reducerende uiteinden.
\textbf{6.} Eén.
\textbf{7.} $12\times 1.9 = \qty{23}{nm}$ straal, \qty{46}{nm}
doorsnede: twee keer een ribosoom.
\textbf{8.} $100/\num{1.43e-17} = \num{7.0e18}$ deeltjes.
\textbf{9.} $100/162 = \qty{0.62}{mol}$ residuen, oftewel \qty{111}{g} vrije
glucose (het water van de \omterm{prop:b1:water-small-molecules:families}{hydrolyse} levert het verschil).
\textbf{10.} \qty{400}{g}, ruim een kwart van de massa van de lever.
\textbf{11.} \qty{10}{g/h} $= \qty{2.78}{mg/s} = \qty{1.54e-5}{mol/s} =
\num{9.3e18}$ moleculen per seconde.
\textbf{12.} Ongeveer \qty{11}{h} (111 g bij 10 g/h) --- een nacht. Vooral aan
die van de hersenen, die \qty{5}{g/h} opnemen.
\textbf{13.} $2048\times\num{7e18} = \num{1.4e22}$ uiteinden.
\textbf{14.} $\num{1.4e22}\times 20 = \num{2.9e23}$ per seconde $=
\qty{0.48}{mol/s} = \qty{86}{g/s}$: dertigduizend keer de werkelijke snelheid. De
uiteinden zijn nooit beperkend; het enzym wel.
\textbf{15.} $\num{9.3e18}/20 = \num{4.6e17}$ moleculen --- ongeveer
\qty{75}{mg} enzym, een klein deel van het levereiwit.
\textbf{16.} Ketens van \num{53000} residuen: nog steeds $\num{7e18}$ uiteinden
(één per keten), \num{2048} keer minder; maximale snelheid \qty{42}{mg/s}, nog
altijd vijftien keer de werkelijke snelheid. (Bij langere amyloseketens krimpt de
marge; de echte grens van een onvertakte voorraad is haar onoplosbaarheid en
kristallisatie, die de uiteinden buiten bereik brengen.)
\textbf{17.} Zodra de bloedglucose daalt, zetten hormonen het fosforylase binnen
seconden aan; het enzym vindt per deeltje duizenden uiteinden die al aan het
oppervlak liggen, zodat de afgifte meteen honderdvoudig kan stijgen, zonder te
wachten tot ketens zijn afgewikkeld of deeltjes zijn opgelost.
\textbf{18.} $0.55/1.0 = \qty{0.55}{osmol/L}$ (\qty{100}{g}
$= \qty{0.55}{mol}$ vrije glucose).
\textbf{19.} Bijna het dubbele van de osmolariteit van het \omterm{def:b1:eukaryotic-cell:organelle}{cytosol}: er zou water
uit het bloed naar binnen stromen, de \omterm{def:b1:cell-unit-of-life:cell}{cellen} zouden tot bijna drie keer hun volume
opzwellen en barsten.
\textbf{20.} $\num{7e18}/\num{6e23} = \qty{1.2e-5}{mol}$ in \qty{1}{L}:
\qty{12}{\micro osmol/L}, veertigduizend keer minder.
\textbf{21.} $\Psi_s = -2.58\times 0.55 = \qty{-1.4}{MPa}$: een wand zou
\qty{1.4}{MPa} moeten dragen, veertien atmosfeer.
\textbf{22.} Spierglycogeen is de eigen brandstof van de spier voor de
samentrekking, die als glucose-6-fosfaat rechtstreeks de glycolyse in wordt
gebracht; de spier mist het enzym dat glucose van haar fosfaat bevrijdt, zodat de
voorraad de vezel niet kan verlaten.
\textbf{23.} Een aardappel mobiliseert haar \omterm{def:b1:carbohydrates:polysaccharide}{zetmeel} over weken van uitlopen; een
lever moet binnen minuten reageren. Grote dichte korrels met weinig uiteinden aan
het oppervlak passen bij een trage, gelijkmatige afgifte; kleine sterk vertakte
deeltjes met duizenden uiteinden bij een snelle.
\textbf{24.} $0.62/(4\times 3600) = \qty{4.3e-5}{mol/s} = \num{2.6e19}$
residuen per seconde; in totaal \qty{1.24}{mol} \omterm{def:b1:water-small-molecules:atp}{ATP}, $\num{5.2e19}$ per seconde.
\textbf{25.} Afgifte: $\num{9.3e18}$ glucosemoleculen per seconde
(\qty{10}{g/h}); mogelijk gemaakt door \num{1.4e22} niet-reducerende uiteinden,
\num{2048} per deeltje; het polymeer bespaart de \omterm{def:b1:cell-unit-of-life:cell}{cellen} \qty{0.55}{osmol/L},
bijna het dubbele van hun eigen osmolariteit.
\end{solution}
